\begin{scanpage}{0238}{220}
{\cyrillicfont\small
\noindent26. Meixner J. und Schäfke F. W. Mathieusche Funktionen und Sphäroidfunktionen. Springer-Verlag, 1954.\par
27. Ландсберг Г. С. Оптика. ГИТТЛ, 1954.\par
28. Braunbek W. Neue Näherungsmethode für die Beugung am ebenen Schirm. Zeits. für Physik, 127, 4, 381--390, 1950.\par
29. Braunbek W. Zur Beugung an der Kreisscheibe. Zeits. für Phys., 127, 4, 405--415, 1950.\par
30. Braunbek W. Zur Beugung an der kreisförmigen Öffnung. Zeits. für Physik 138, 1, 80--88, 1954.\par
31. Braunbek W. Zur Beugung an Öffnungen in nichtebenen Schirmen. Zeits. für Physik, 156, 1, 65, 1959.\par
32. Frahn W. E. Beugung elektromagnetischer Wellen in Braunbekscher Näherung. I. Zeits. für Physik, 156, 1, 78, 1959.\par
33. Frahn W. E. Beugung elektromagnetischer Wellen in Braunbekscher Näherung. II. Zeits. für Phys., 156, № 2, 99, 1959.\par
34. Белкина М. Г. Дифракция электромагнитных волн на диске. В сб. «Дифракция электромагнитных волн на некоторых телах вращения». Изд-во «Советское радио», 1957, (стр. 148--174).\par
35. Иванов В. И. Коротковолновая асимптотика дифракционного поля в тени идеально проводящего параболического цилиндра. «Радиотехника и электроника», 5, № 3, 393--402, 1960.\par
36. Иванов В. И. Дифракция коротких плоских электромагнитных волн на гладком выпуклом цилиндре при наклонном падении. «Радиотехника и электроника», 5, № 3, 524--528, 1960.\par
37. Горяинов А. С. Асимптотическое решение задачи о дифракции плоской электромагнитной волны на проводящем цилиндре. «Радиотехника и Электроника», 3, № 5, 603--614, 1958.\par
38. Felsen L. B. Backscattering from wide-angle and narrow-angle cones. Journ. Appl. Phys. 26, № 3, 138--151, 1955.\par
39. Siegel K. M., Crispin J. W. and Schensted C. E. Electromagnetic and acoustical scattering from a semi-infinite cone. Journ. Appl. Phys., 26, № 3, 309--313, 1955.\par
40. Горяинов А. С. Дифракция плоской электромагнитной волны, распространяющейся вдоль оси конуса. «Радиотехника и электроника», 6, № 1, 47--57, 1961.\par
41. Siegel K. M., Goodrich R. F., Weston V. H. Comments on far field scattering from bodies of revolution. Appl. Sci. Res., Section B, volume 8, p. 8, 1959.\par
42. Keller J. B. Diffraction by a convex cylinder. Trans. IRE, AP-4, № 3, 312--321, 1956.\par
43. Keller J. B. Diffraction by an aperture. I. Journal of Appl. Phys. 28, № 4, 426--444, 1957. Diffraction by an aperture. II. Journ. Appl. Phys., 28, № 5, 570--579, 1957.\par
44. Keller J. B. Backscattering from a finite cone. Trans. IRE, AP-8, № 2, 175--182, 1960. Backscattering from a finite cone---comparison of theory and experiment. Trans. IRE, AP-9, № 4, 411, 1961.\par
45. Schensted C. E. Electromagnetic and acoustic scattering by a semi-infinite body of revolution. Journ. Appl. Phys., 26, № 3, 306--308, 1955.\par
46. Clemmow P. C. Edge currents in diffraction Theory. Transactions IRE. Antennas and Propagation, AP-4, № 3, 282--287, 1956.\par
47. Millar R. F. An approximate theory of the diffraction of an electromagnetic wave by an aperture in a plane screen. Proc. Inst. Electrical Engrs., Part C, 103, № 3, 177, 1956.\par
48. Millar R. F. The diffraction of an electromagnetic wave by a circular aperture. Proc. Inst. Electrical Engrs., Part C, 104, № 5, 1957.\par
49. Millar R. F. The diffraction of an electromagnetic wave by a large aperture. Proc. Inst. Electrical Engrs., Part C, 104, № 6, 1957.\par
50. Handbuch der Physik. Bd 25/1. Kristalloptik. Beugung. Berlin, Springer, 1961. H. Hönl, A. W. Maue, K. Westpfahl. Theorie der Beugung (S. 218--592).\par
51. Karp S. N. and Russek A. Diffraction by a wide slit. Journal of Applied Physics, 27, № 8, 886--891, 1956.\par
52. Buchal R. N. and Keller J. B. Boundary Layer Problems in Diffraction Theory. Commun. Pure and Appl. Mathem., 13, № 1, 85--114, 1960.\par}
\end{scanpage}
